\documentclass[11pt,a4paper]{article}
\usepackage[utf8]{inputenc}
\usepackage[T1]{fontenc}
\usepackage{lmodern}
\usepackage{xcolor}
\renewcommand{\contentsname}{Spis treści}
\renewcommand{\tablename}{Tabela}
\hyphenpenalty=3000 \tolerance=2000 \emergencystretch=2em
\usepackage[margin=2.1cm]{geometry}
\usepackage{booktabs}
\usepackage{tabularx}
\usepackage{longtable}
\usepackage{enumitem}
\usepackage[colorlinks=true,linkcolor=blue!50!black,urlcolor=blue!50!black]{hyperref}
\usepackage{titlesec}
\usepackage{parskip}
\titleformat{\section}{\large\bfseries}{\thesection.}{0.6em}{}
\titleformat{\subsection}{\normalsize\bfseries}{\thesubsection.}{0.6em}{}
\titleformat{\subsubsection}{\small\bfseries}{\thesubsubsection.}{0.6em}{}
\setlist{nosep,leftmargin=1.4em}
\newcommand{\ohm}{$\Omega$}

\title{\textbf{AMPERE POINT --- Wielofunkcyjne urządzenie pomiarowe (custom device)}\\[4pt]
\large Koncepcja --- wersja 6: PEŁNY PLAN REALIZACJI PROJEKTU\\
zautomatyzowana stacja testowa EVSE z budżetowym elektrycznym magazynem energii\\[2pt]
\normalsize (dokument samodzielny --- nie wymaga znajomości wersji 1--5)}
\author{Opracowanie wewnętrzne AMPERE POINT}
\date{9 lipca 2026}

\begin{document}
\maketitle
\vspace{-1.4em}

\section*{Streszczenie wykonawcze}
Budujemy \textbf{stację testową ładowarek AC (EVSE)}, która wobec badanej ładowarki udaje samochód
elektryczny: symuluje stany złącza wg IEC~61851-1, pobiera zadany prąd (1F do 32~A, 3F do 16~A)
i rejestruje wyniki --- bez auta testowego. Energia testów \textbf{nie jest tracona}: trafia przez
używane prostowniki telekomunikacyjne (6$\times$ Eltek Flatpack2 48/2000~HE, sterowane po CAN) do
\textbf{baterii 48~V składanej samodzielnie} (16 ogniw LFP 105~Ah, 5,4~kWh, rozbudowa do 10,8~kWh),
a z niej --- przez falownik zaadaptowany z używanego UPS-a --- zasila wydzielone obwody warsztatu.
System składa się z \textbf{modułu A} (przenośna jednostka sterująco-pomiarowa z testerem PM701E
obracanym serwami) i \textbf{szafy B} (magazyn energii). Bezpieczeństwo opiera się na sprzętowym
łańcuchu 24~V, którego procesor nie może zmostkować. Budżet: \textbf{$\sim$9,4--17,4~tys.~zł}
(w tym bateria --- koszt nieredukowalny); praca własna: $\sim$6--10 dni roboczych. Dokumenty
wykonawcze: schemat elektryczny v5 (EPLAN, arkusze E1--E3), objaśnienie z przewodnikiem
uruchomienia, wizualizacja 3D v3.

\tableofcontents

\section{Cel, kontekst i wymagania}
\subsection{Kontekst}
Zakład AMPERE POINT importuje, testuje, sprzedaje i serwisuje ładowarki AC do aut elektrycznych
(wallboxy 11/22~kW serii PRIME/HM, przenośne Q/P/B; Type~2; firmware producenta zamknięty).
Potrzebne jest stanowisko, które pozwoli: (a) wykonywać \textbf{obowiązkowe pomiary elektryczne}
stacji ładowania, (b) prowadzić \textbf{testy funkcjonalne pod realnym obciążeniem} (kontrola
importu, serwis, regresje po nowym firmware) --- obecnie robione autem testowym, które ma przestać
być potrzebne.
\subsection{Pięć obowiązkowych pomiarów (podstawa: ,,Przewodnik pomiarów stacji ładowania'', W-wa 2024)}
\begin{enumerate}
\item \textbf{Ciągłość przewodów ochronnych} --- prąd $\ge$200~mA, kryterium niskiej rezystancji.
\item \textbf{Rezystancja izolacji} --- 250/500/1000~V~DC, $\ge$1~M\ohm{} (250~V dla obwodów
z elektroniką; żyły czynne można zewrzeć i mierzyć względem PE).
\item \textbf{Rezystancja uziemienia} --- metoda techniczna 3-elektrodowa (E/S/H), gdy stacja ma
własny uziom.
\item \textbf{Sprawdzenie RCD} --- typ, prąd i czas zadziałania, napięcie dotykowe U$_B$;
dodatkowo RDC-DD 6~mA~DC (IEC~62955).
\item \textbf{Impedancja pętli zwarcia Z$_s$} --- warunek Z$_s\cdot$I$_a\le$U$_o$.
\end{enumerate}
Pomiary 1, 2, 4, 5 wymagają wprowadzenia ładowarki w stan ,,ładuje'' i dostępu do żył --- stąd
w systemie tester-adapter PM701E i własne gniazdo pojazdowe.
\subsection{Minimum teorii: interfejs ładowarka--pojazd (IEC 61851-1)}
Ładowarka AC to ,,inteligentny stycznik'' --- prostownik jest w aucie (OBC). Jedyna rozmowa odbywa
się linią \textbf{CP} (Control Pilot): ładowarka podaje $\pm$12~V; pojazd rezystorami ustawia stany:
\textbf{A} $+12$~V (brak auta), \textbf{B} $+9$~V (podłączony), \textbf{C} $+6$~V (ładuj --- ładowarka
zamyka stycznik), \textbf{D} $+3$~V, \textbf{E/F} 0/$-12$~V (błąd). PWM 1~kHz na CP deklaruje prąd:
I[A] $=$ wypełnienie[\%]$\times$0,6 (np. 53\% $\approx$ 32~A). Linia \textbf{PP} koduje rezystancją
obciążalność kabla (1,5~k\ohm=13~A, 680~\ohm=20~A, 220~\ohm=32~A, 100~\ohm=63~A). Zabezpieczenie
\textbf{RDC-DD} musi wykryć gładki prąd stały upływu $\ge$6~mA (RCD typu~A go nie widzi).
\subsection{Wymagania funkcjonalne stacji}
\begin{itemize}
\item symulacja stanów A/B/C/D/E z weryfikacją; deklaracje prądu 13/20/32/63~A (pokrętła PM701E);
\item pobór prądu: 1F 6--32~A oraz 3F do 16~A/fazę (11~kW), kroki $\le$25\% deklaracji, rampy,
soak 15--30~min z rejestracją; pobór sinusoidalny (PF$\approx$1) jak prawdziwy OBC;
\item testy zabezpieczeń: RDC-DD 6~mA~DC automatycznie (rampa, pomiar prądu i czasu zadziałania),
RCD typ~A protokołowo miernikiem wzorcowanym; pętla, izolacja, ciągłość --- miernikiem wzorcowanym
w sekwencji prowadzonej;
\item termowizja wtyku podczas soak ($\Delta$T, hot-spoty); pomiar U/I/P/cos$\varphi$ per faza;
\item magazynowanie energii testów i jej użycie w zakładzie; brak konieczności bieżącego zużywania;
\item rejestracja wyników (SD + MQTT$\to$Home Assistant), protokół PDF na PC;
\item bezpieczeństwo: rozłączanie mocy sprzętowe, niezależne od firmware.
\end{itemize}

\section{Zasady projektowe}
\begin{enumerate}
\item \textbf{Spójność do poziomu złączy i sygnałów} --- architektura przenosi się 1:1 na schemat.
\item \textbf{Bezpieczeństwo sprzętowe, nie programowe} --- szeregowy łańcuch styków 24~V zasila
cewkę stycznika mocy; firmware może tylko zezwalać/rozłączać (styk -K10) i musi dowodzić życia
(watchdog -A2); żadnego członu nie da się zmostkować programowo.
\item \textbf{Praktyczna wykonalność} --- komponenty z rynku (w tym wtórnego), budowa etapowa,
każdy etap kończy się czymś użytecznym.
\item \textbf{Wiarygodność protokołów} --- pomiary urzędowe wykonują przyrządy wzorcowane
(Uni-T UT595, UT522) prowadzone przez stację; bloki własne robią to, czego przyrządy nie umieją.
\item \textbf{Tanio przez pracę własną} --- używane Flatpack2 zamiast nowych prostowników, pakiet
bateryjny składany samodzielnie, falownik z używanego UPS-a, szyny DC z płaskownika, ręczne pole
przełączeń zamiast stycznika konfiguracyjnego.
\end{enumerate}

\section{Architektura systemu}
\subsection{Dwa moduły i przepływ energii}
\begin{center}\small
DUT (ładowarka, zasilana z instalacji) $\to$ wtyk Type~2 $\to$ \textbf{MODUŁ A} [gniazdo -X1 $\to$
rozłącznik -Q0 $\to$ bezpieczniki -F1..3 $\to$ czujnik różnicowy -B1 $\to$ przekładniki -T1..3 $\to$
stycznik -K1 $\to$ gniazdo CEE -X2] $\to$ przewód 5$\times$6~mm$^2$ $\to$ \textbf{SZAFA B}
[styczniki -K11..13 $\to$ 6$\times$ Flatpack2 $\to$ szyna DC 48~V (-F13, -KB, -Q1) $\to$ pakiet
-GB1 $\to$ falownik (UPS) $\to$ wydzielone obwody wyspy].
\end{center}
Równolegle z -X1 wychodzi \textbf{odczep sterujący -X3} (7 żył: L1/L2/L3/N/PE/CP/PP, $\le$1~m) do
testera \textbf{PM701E}, którego pokrętła obracają serwa. Prąd obciążenia NIE płynie przez PM701E
(adaptery tej klasy nie przenoszą ciągłych 32~A) --- to decyzja topologiczna Z1. Moduł A jest
przenośny (pomiary terenowe robi sam, z UT595/UT522); szafa B jest warsztatowa.
\subsection{Dlaczego prostownik+bateria to najlepsza symulacja auta}
Pokładowa ładowarka auta (OBC) to prostownik z aktywnym PFC ładujący baterię. Flatpack2 to
dokładnie taki prostownik: pobór sinusoidalny, PF$>$0,99, THD niskie, prąd zadawany cyfrowo ---
zachowanie widziane przez DUT jest nieodróżnialne od auta, a regulacja jest płynna (rampy,
asymetria faz), czego bank rezystorów nie potrafi.
\subsection{Konfiguracja mocy}
\begin{longtable}{@{}p{0.30\textwidth}p{0.40\textwidth}p{0.22\textwidth}@{}}
\toprule
\textbf{Tryb} & \textbf{Przydział Flatpack2 (2 kW/szt.)} & \textbf{Pokrycie}\\
\midrule
3F symetrycznie do 16 A (11 kW) & po 2 szt. na fazę (4 kW $>$ 3,7 kW/fazę) & pełne\\
1F do 32 A (7,4 kW) & 4 szt. na L1: 2 własne $+$ 2 przepięte mostkami na polu -X6 & pełne\\
margines 3$\times$32 A & wymagałby 12 szt. & opcja przyszła\\
\bottomrule
\end{longtable}
Pole przełączeń \textbf{-X6} to ręczne mostki (zamiast stycznika) --- oszczędność; poprawność
konfiguracji sekwencer weryfikuje pomiarem prądów per faza (-U2) przy 25\% mocy, zanim pozwoli
na pełny test.

\section{Moduł A --- bloki funkcjonalne}
\subsection{B1: wejście i tor mocy}
\textbf{-X1} gniazdo pojazdowe Type~2 32~A~3F z blokadą elektromagnesową \textbf{-Y1} (rygluje
wtyk w stanie C). \textbf{-X3} odczep 7-żyłowy do PM701E; w linii CP zestyk przekaźnika
\textbf{-K4}: E-STOP lub zanik 24~V przerywa CP $\Rightarrow$ DUT widzi odłączenie pojazdu i sam
otwiera swój stycznik (magistrala gaśnie); na listwie mostek serwisowy. \textbf{-Q0} rozłącznik
4P (uwaga: nie odłącza odczepu --- tabliczka przy gniazdach bananowych). \textbf{-F1..-F3} gG
32~A na fazach (N bez wkładki, rozłączany 4-biegunowo). \textbf{-B1} czujnik różnicowy AC/DC
(fluxgate, np. LEM CDSR: okno L1+L2+L3+N) --- monitor RCM własnych torów i weryfikacja generatora
upływu; SPI do -A1. \textbf{-T1..-T3} przekładniki 40~A/1~V do meteringu. \textbf{-K1} stycznik
główny 4P 40~A AC-1 --- jedyny aparat podający moc do szafy; cewka zasilana wyłącznie z łańcucha
bezpieczeństwa. \textbf{-X2} gniazdo CEE 32~A~5p. Przekroje toru 6~mm$^2$; PE nieprzerwane.
\subsection{B2: symulacja stanów --- PM701E + serwa}
PM701E (posiadany): pokrętło stanu BB/C/A/B/D i prądu NC/13/20/32/63~A, gniazda bananowe
L1/N/PE/L2/L3 (punkt przyłączenia UT595), wyjście SIGNAL~OUTPUT (CP). Dwa serwa \textbf{-M1/-M2}
(DS3218, 20~kg$\cdot$cm) na ramie nieinwazyjnej nad pokrętłami (adapter drukowany + sprzęgło
Oldham); kalibracja kąt$\leftrightarrow$pozycja (5+5 pozycji); \textbf{każda zmiana potwierdzana
odczytem CP} --- serwo nie jest zaufane. Przyciski PM701E (Touch, PE Pre-Test, CP~Error, PE~Error)
w etapie 1 ręczne.
\subsection{B4: metering}
\textbf{-U2} ATM90E36A (front-end klasy licznikowej, SPI): 3 napięcia (dzielniki odłączane
przekaźnikami separującymi \textbf{-K31..-K33} na czas prób izolacji) + 3 prądy (-T1..3);
U/I/P/Q/cos$\varphi$/f/THD per faza; rola: weryfikacja faz po wejściu w stan C, kontrola
zadane-vs-pobrane (niezależnie od CAN), rejestracja soak, weryfikacja konfiguracji -X6.
Optotransoptory obecności L1/L2/L3 --- pomiar czasów zadziałania zabezpieczeń DUT (znacznik ms).
\subsection{B5: generator upływu (test RDC-DD/RCD)}
\textbf{-A3}: źródło prądowe 0--20~mA (DAC $\to$ wzmacniacz $\to$ MOSFET; pomiar bocznikiem przez
wzmacniacz izolowany AMC1301), przyłącze L1 \emph{za} oknem -B1 (przez -K30) $\to$ PE. Rampa
(0,1~mA/s przy progu) $\to$ prąd zadziałania RDC-DD ($\le$6~mA $=$ PASS) i czas odcięcia (z opto
faz). Tryb AC do testów przesiewowych RCD typu~A; pomiar protokołowy RCD wykonuje UT595.
\subsection{B6: akwizycja CP/PP}
Z SIGNAL~OUTPUT PM701E: dzielnik ($\pm$12~V$\to$0--3,3~V) + bufor + ograniczniki; ADC mierzy
poziom stanu, komparator $\to$ wejście przechwytujące mierzy wypełnienie PWM 1~kHz (deklaracja
prądu). PP: drugi kanał ADC (kod kabla). Funkcje: weryfikacja serw, protokół, wykrywanie błędów.
\subsection{B7: termika}
\textbf{-B2} MLX90640 (matryca IR 32$\times$24) nad zatoką -X1: $\Delta$T wtyku i hot-spoty
podczas soak; \textbf{-B3/-B4} MLX90614 (pirometry punktowe, drugi adres!); DS18B20 (1-Wire):
szyny, -K1, radiatory Flatpack2, ogniwa pakietu, otoczenie (odniesienie).
\subsection{B8: rdzeń przetwarzający -A1}
ESP32-S3-WROOM (2 rdzenie, WiFi, 45 GPIO). Interfejsy: \textbf{SPI} (-B1, -U2, microSD),
\textbf{I2C} (MLX-y, ekspander MCP23017 $\to$ ULN2803 $\to$ cewki, RTC DS3231, DAC -A3 za
izolatorem), \textbf{1-Wire} (DS18B20), \textbf{UART} (HMI DWIN 7$''$ w pokrywie), \textbf{CAN/TWAI
przez transceiver izolowany} (Flatpack2 $\times$6 + BMS), \textbf{ADC} (CP, PP, bocznik -A3),
\textbf{PWM} (serwa), \textbf{DO przez opto} (DO1 zezwolenie -K10, DO2 impulsy watchdog -A2,
DO3 blokada -Y1), \textbf{DI przez opto} (stan członów łańcucha, opto faz). Zasilanie: -G1
(230~V z własnego przewodu Schuko -X4, przez -F4): 24~V (łańcuch, cewki), 5/3,3~V (logika, serwa)
--- nigdy z badanej ładowarki.
\subsection{B9: łańcuch bezpieczeństwa 24 V (arkusz E2 schematu)}
Szeregowo od +24~V: \textbf{-S0} E-STOP (drugi zestyk zwalnia -K4/CP) $\to$ \textbf{-F11.1..8}
termiki KSD301 na radiatorach Flatpack2 i falownika (NC) $\to$ \textbf{-S1} krańcówka pokrywy
strefy mocy $\to$ \textbf{-K2} potwierdzenie prądowe wentylatorów $\to$ \textbf{-K10} zezwolenie
MCU $\to$ \textbf{-A2} watchdog impulsowy (zanik impulsów z -A1 $=$ rozpad) $\to$ [\textbf{-S2}
START $\parallel$ samopodtrzymanie] $\to$ cewka \textbf{-RB} (przekaźnik bezpieczeństwa, styki
wymuszone); drugi szczebel: styk -RB $\to$ cewka -K1. Każda awaria/przerwa $=$ -K1 otwiera moc;
powrót wyłącznie przyciskiem START. Strona DC: -F13.1/2 (mega), -Q1 (rozłącznik serwisowy),
\textbf{-KB} kontaktor sterowany przez BMS (odcięcie baterii niezależne od wszystkiego).
\subsection{B10: panel i gniazda pomocnicze}
E-STOP, START, napęd -Q0, kolumna LED (czerwony: napięcie na magistrali; żółty: test; zielony:
gotowość), gniazda bananowe pomiarowe L1/L2/L3/N/PE (tabliczka: mogą być pod napięciem w stanie C),
zaciski E/S/H (uziom, UT522), HMI w pokrywie, przepust Schuko, USB/SD serwisowo.

\section{Szafa B --- budżetowy magazyn energii}
\subsection{Prostowniki: 6$\times$ Eltek Flatpack2 48/2000 HE (używane)}
Moduły telekomunikacyjne z demontaży ($\sim$150--300~zł/szt.): 2~kW, wejście 185--275~V~AC,
wyjście 43,5--57,5~V~DC, sprawność 96,5\%, PF$>$0,99, chłodzenie własnym wentylatorem, złącze
krawędziowe (kupić z gniazdami lub półką 2--4-slotową). \textbf{Sterowanie CAN} (125~kb/s):
zadawanie napięcia/limitu prądu i odczyt telemetrii (protokół opisany społecznościowo ---
weryfikacja w etapie 0; fallback: Huawei R4850G2). Przydział faz przez styczniki \textbf{-K11..-K13}
(cewki 24~V w torze zezwolenia) i pole mostków \textbf{-X6} dla 1F 32~A. Zakup: 7 szt. (1 zapas),
sprawdzenie każdej na stole (rozdz.~\ref{sec:plan}).
\subsection{Szyna DC i zabezpieczenia}
Płaskownik Cu 20$\times$3~mm na izolatorach; przewody 50~mm$^2$ (przy 11~kW $\to\sim$230~A);
\textbf{-F13.1/-F13.2} bezpieczniki mega 150--200~A przy pakiecie; \textbf{-Q1} rozłącznik DC
serwisowy; \textbf{-KB} kontaktor główny sterowany stykiem BMS; wszystkie połączenia na M8
z kontrolą momentu (wkrętak dynamometryczny --- posiadany).
\subsection{Pakiet bateryjny -GB1: 16S LFP 105 Ah (DIY) --- przepis}
\label{sec:pakiet}
\textbf{Wymiarowanie:} soak 30~min@11~kW $=$ 5,5~kWh. Start: 1 gałąź 16S (51,2~V, 5,4~kWh) ---
pełny soak 1F (3,7~kWh) i 3F przy SOC startowym $\le$25\% lub z równoległym rozładowaniem przez
falownik; docelowo 2P (10,8~kWh): pełny soak 3F przy SOC $\le$40\%. Ogniwa: LFP 100--105~Ah
klasy~B (np. EVE LF105) z pomiarami od sprzedawcy.
\begin{enumerate}
\item \textbf{Selekcja:} zmierz napięcia (rozrzut $\le$30~mV) i rezystancję wewnętrzną; paruj.
\item \textbf{Top-balance:} wszystkie ogniwa równolegle, ładowanie CV 3,55--3,60~V do prądu
$\approx$0 (12--48~h). Źródło: zasilacz lab. albo jeden Flatpack2 ustawiony na 3,6~V$\times$...
(po zbudowaniu 16S: 56,8~V dla balansu górnego całej gałęzi).
\item \textbf{Kompresja:} płyty czołowe + 4 pręty M8, docisk wg noty ($\sim$300~kgf).
\item \textbf{Połączenia:} busbary fabryczne/płaskownik, momenty M6 4--6~N$\cdot$m; pomiar spadków
na łączach przy prądzie próbnym.
\item \textbf{BMS} 16S 150--200~A z balansem aktywnym (klasy JK): przewody balansera od minusa
kolejno; progi 2,80/3,55~V, ładowanie 0--45~$^\circ$C, 2$\times$NTC; wyjście BMS $\to$ cewka -KB.
\item \textbf{Zabezpieczenie gałęzi:} mega 150--200~A przy $+$, przewody 50~mm$^2$.
\item \textbf{Pierwsze ładowanie nadzorowane} (10--20~A z jednego FP2), obserwacja ogniw w aplikacji.
\item \textbf{Pomiar pojemności:} rozładowanie znanym obciążeniem przez falownik, zapis do metryki.
\item \textbf{Zabudowa:} dolny przedział szafy, przegroda od AC, wentylacja grawitacyjna, DS18B20
na ogniwach, etykieta (pojemność, data, progi).
\end{enumerate}
\subsection{Falownik W2: adaptacja używanego UPS on-line --- przepis}
\label{sec:ups}
\textbf{Kryteria zakupu} (500--1500~zł): UPS on-line (podwójna konwersja) 5--6~kVA w wieży,
\textbf{szyna baterii 48~V} (typowe w tej klasie --- potwierdzić w dokumentacji konkretnego
modelu przed zakupem!), sprawny tor falownika, dostępne złącze baterii zewnętrznej.
\begin{enumerate}
\item Test wstępny UPS-a z jego akumulatorami / zasilaczem: wyjście 230~V, przebieg (mikrooscyloskop).
\item Wyprowadź tor baterii: złącze zewnętrzne $\to$ przewody 25--35~mm$^2$ $\to$ przez -F13/-Q1
do szyny -GB1 (napięcie pakietu 44--57~V mieści się w oknie UPS 48~V).
\item Wyłącz/ogranicz ładowarkę wewnętrzną UPS (ładowaniem zarządzają Flatpack2 i BMS): najlepiej
programowo (menu/serwis) lub odłączeniem jej toru --- zależnie od modelu.
\item Skonfiguruj progi rozładowania UPS powyżej progu BMS (np. 44~V), alarmy dźwiękowe wycisz.
\item Wyjście UPS $\to$ rozdzielniczka wyspy: B16 + RCD 30~mA typ~A + gniazda; \textbf{punkt N--PE
wyspy} przy źródle (inaczej RCD nie zadziała); obwody oznakowane, zasilanie wybierane
przełącznikiem \textbf{sieć--0--wyspa} (nigdy równolegle z siecią!).
\item Testy: praca z pakietu 1/2/4~kW (grzałka/czajnik/narzędzia), zapady przy udarze silnikowym,
1~h soak z pomiarem temperatur; pomiar RCD wyspy (UT595).
\end{enumerate}
Alternatywa \textbf{W3} (gdy UPS się nie trafi): własny falownik wyspowy LF 48/230~V, 3--5~kW ---
kompletny przepis znajduje się w koncepcji v3, rozdz.~6 (EG8010/EGS002 + most H + toroid);
BOM 1,35--2,5~tys.~zł. Upgrade przyszłościowy \textbf{W1}: fabryczny falownik hybrydowy ESS ---
wtedy dochodzi oddawanie do sieci (zgłoszenie mikroinstalacji) bez zmian reszty szafy.
\subsection{Wentylacja i złącza}
Straty szafy 0,4--0,9~kW (sprawności FP2/UPS) --- 2 wentylatory szafowe (-M3/-M4, zasilanie z -X4
przez -F5, załączane -K3, potwierdzane przekaźnikiem prądowym -K2 w łańcuchu). Złącza: CEE (moc),
\textbf{-X5} (sterowanie: cewki K11--K13, pętla termików, wentylatory), CAN (skrętka z transceiverem),
DC do UPS-a.

\section{Zarządzanie energią (bez zrzutu cieplnego)}
Tryby: \textbf{TEST} (FP2 ładują pakiet prądem zadanym przez sekwencer; falownik może równolegle
zasilać wyspę), \textbf{MAGAZYN/UPS} (poza testami wyspa działa z pakietu), \textbf{SERWIS}
(-Q1 otwarty). \textbf{Bramka SOC} przed soakiem: 3F pełny --- SOC$\le$40\% (2P) / $\le$25\% (1P);
1F --- $\le$60\%; przy przekroczeniu HMI: ,,rozładuj magazyn'' (włącz odbiory wyspy). Telemetria:
BMS (SOC, napięcia ogniw, temperatury) po RS485/CAN; FP2 (U/I/T) po CAN; -U2 liczy energię pobraną
z DUT --- bilans i sprawność w raporcie. Plan C (nieobowiązkowy, $\sim$100~zł): rezystor 2~kW
załączany ręcznie do szybkiego zrzutu.

\section{Bezpieczeństwo --- zestawienie}
\begin{itemize}
\item \textbf{AC}: łańcuch E2 (rozdz. 4.8), -K4 gasi CP przy E-STOP, gG na fazach, RCM (-B1)
zdejmuje zezwolenie, blokada wtyku -Y1, tabliczki przy gniazdach bananowych.
\item \textbf{DC}: mega przy pakiecie, -Q1, -KB od BMS, przekroje/momenty, przedział bateryjny
z przegrodą; LFP (chemia stabilna) zamiast NMC.
\item \textbf{Wyspa}: brak pracy równoległej z siecią (przełącznik z przerwą), N--PE przy źródle,
RCD 30~mA, oznakowanie obwodów.
\item \textbf{Formalnie}: użytek własny zakładu; wykonanie i odbiór --- osoba z uprawnieniami
SEP E/D; pomiary protokołowe przyrządami z aktualnym świadectwem; zgodność ogólna: IEC~61010
(kategorie pomiarowe), IEC~61851 (interfejs), PN-HD~60364 (instalacja, wyspa).
\end{itemize}

\section{Sekwencja pomiarowa (pełny przebieg testu ładowarki)}
\begin{enumerate}
\item \textbf{Self-test:} odczyt członów łańcucha, próba wentylatorów (-K3$\to$-K2), komunikacja
(CDSR, metering, BMS, FP2, HMI), serwa do pozycji spoczynkowych (BB/NC) z potwierdzeniem CP,
\textbf{bramka SOC}.
\item \textbf{Tryb izolacji:} -K30..-K33 otwarte (elektronika własna odłączona od magistrali);
HMI prowadzi: UT595 na gniazdach bananowych PM701E, próba 250~V (elektronika DUT!), żyły czynne
zwarte względem PE; wynik $\ge$1~M\ohm{} wpisany $=$ bramka.
\item \textbf{Ciągłość PE:} UT595 $\ge$200~mA przez wyjście PE; wpis $=$ bramka.
\item \textbf{Stan B:} serwo stanu $\to$ B, CP musi pokazać $+9$~V; pre-testy PM701E.
\item \textbf{Stan C:} serwo $\to$ C, CP $+6$~V; DUT zamyka stycznik; -U2 weryfikuje napięcia
i kolejność faz; -Y1 rygluje wtyk.
\item \textbf{Zabezpieczenia:} RDC-DD automatycznie (-A3: rampa DC, prąd zadziałania $\le$6~mA,
czas z opto faz); RCD typ~A protokołowo (UT595: $\frac{1}{2}$I$_{\Delta n}$, I$_{\Delta n}$,
5I$_{\Delta n}$, U$_B$); po każdym zadziałaniu prowadzony reset DUT i powrót do C.
\item \textbf{Pętla Z$_s$:} UT595 na gniazdach bananowych; wpis.
\item \textbf{Obciążenie/soak:} serwo prądu $\to$ deklaracja; kontrola konfiguracji -X6 przy 25\%
(zgodność faz z -U2); kroki 25/50/75/100\% deklaracji zadawane FP2 po CAN (zawsze $\le$ deklaracji
CP); soak 15--30~min: log U/I/P per faza, $\Delta$T wtyku (MLX90640), temperatury pakietu;
kryteria pass/fail (spadki, asymetria, $\Delta$T, stabilność).
\item \textbf{Testy błędów:} przyciski CP~Error ,,E'' / PE~Error na PM701E (etap 1 ręcznie) ---
czas reakcji DUT z opto faz.
\item \textbf{Uziom} (niezależnie): UT522 na E/S/H.
\item \textbf{Raport:} CSV/JSON na SD + MQTT$\to$HA; protokół PDF na PC; metryka energii
(pobrana/zmagazynowana/sprawność).
\end{enumerate}
Warunki ciągłe: CP zgodne z komendą, RCM $<$ progu, termiki OK, wentylatory OK, BMS OK, E-STOP
zwolniony --- naruszenie $=$ przerwanie (zezwolenie zdjęte, zrzut obciążenia CAN-em, serwa do
spoczynku, log przyczyny).

\section{Firmware -A1 --- specyfikacja funkcji}
\begin{itemize}
\item \textbf{Sekwencer} stanów z bramkami (izolacja $\to$ PE $\to$ B $\to$ C $\to$ testy $\to$
soak) i ścieżkami błędów; tryb serwisowy ręczny (pojedyncze aparaty z HMI, z blokadami).
\item \textbf{Serwa:} tabela kalibracji 5+5 pozycji (EEPROM), najazd, weryfikacja CP, procedura
kalibracyjna z HMI.
\item \textbf{Akwizycja:} CP (poziom+duty), PP, -U2 (200~ms), CDSR (50~ms, progi zależne od fazy
testu), DS18B20 (2~s), bocznik -A3 (1~ms w teście RDC-DD), opto faz (przerwania, znaczniki ms).
\item \textbf{CAN:} sterownik FP2 (napięcie 55,2~V ładowania, limit prądu per moduł, załączanie),
odczyt BMS; watchdog komunikacji (brak ramek 1~s $=$ przerwanie testu).
\item \textbf{Bezpieczeństwo programowe} (uzupełniające, nie zastępuje łańcucha): spójność
CP$\leftrightarrow$suma limitów FP2, bramka SOC, limity temperatur, timeout stanów.
\item \textbf{HMI (DWIN):} ekrany: start/status, prowadzenie pomiarów UT595 z polami wpisu,
przebieg soak (wykres), alarmy, serwis, kalibracja; fizyczny START/E-STOP poza HMI.
\item \textbf{Rejestracja:} SD (CSV: telemetria 1~s; JSON: wynik testu), RTC DS3231, identyfikator
DUT (wpis/QR w przyszłości); \textbf{MQTT} do Home Assistant zakładu (podgląd, archiwum).
\item \textbf{OTA/serwis:} aktualizacja po USB; konfiguracja (progi, moce) w pliku na SD.
\end{itemize}

\section{Obudowy i rozmieszczenie}
\textbf{Moduł A} (rozdzielnica przenośna IP54 $\sim$600$\times$400$\times$250, przezroczysta
pokrywa): strefa mocy za osłoną z krańcówką (lewa), strefa analogowa (środek), logika (prawa);
półka PM701E z ramą serw pod pokrywą; HMI w pokrywie; panel: -X1, -X2, E-STOP, START, napęd -Q0,
LED, banany, E/S/H, przepust -X4. \textbf{Szafa B} ($\sim$600$\times$1200$\times$400 na kółkach):
dół --- pakiet -GB1 za przegrodą; środek --- szyna DC, -F13, -KB, -Q1, BMS; wyżej --- półka
6$\times$FP2; góra --- -K11..13, pole -X6, wentylatory; obok --- wieża UPS (W2) połączona DC;
wyjście na rozdzielniczkę wyspy (B16+RCD) z przełącznikiem sieć--0--wyspa. Szczegóły przestrzenne:
\textbf{wizualizacja 3D v3} (plik HTML w folderze).

\section{Plan realizacji --- etapy, czasy, zakupy}
\label{sec:plan}
\begin{longtable}{@{}p{0.06\textwidth}p{0.46\textwidth}p{0.20\textwidth}p{0.20\textwidth}@{}}
\toprule
\textbf{Etap} & \textbf{Zakres i zadania} & \textbf{Zakupy kluczowe} & \textbf{Czas własny}\\
\midrule
0 & Weryfikacje na stole: (a) serwa na PM701E --- rama, kalibracja, trwałość, potwierdzanie CP;
(b) tor CP (dzielnik+capture) na płytce; (c) \textbf{1$\times$FP2 + CAN z ESP32}: zadawanie prądu,
telemetria; (d) próbka CDSR; (e) stanowisko top-balance & 1$\times$FP2, transceiver CAN, CDSR,
serwa, drobnica & 3--4 dni\\
1 & Moduł A kompletny (tor mocy, łańcuch E2 z -K4, bloki pomiarowe, HMI, firmware rdzeń) +
\textbf{szafa 1F}: 3--4$\times$FP2 na L1, pakiet 16S1P (przepis \ref{sec:pakiet}), UPS W2
(przepis \ref{sec:ups}), wyspa; pełny test ładowarek 1F do 32~A z magazynowaniem & obudowy,
aparaty E1/E2, ogniwa 16 szt.+BMS, UPS, FP2 do 4 szt. & 4--6 dni\\
2 & Rozbudowa 3F: 6$\times$FP2, pole -X6, (opcjonalnie) druga gałąź 2P; testy asymetrii;
aktuatory przycisków PM701E & 2$\times$FP2, 16 ogniw (2P) & 1--2 dni\\
3 & Rozszerzenia: upgrade falownika do W1 (ESS, zgłoszenie), SZR wyspy, front-endy przesiewowe,
wzorcowanie bloków własnych & wg decyzji & ---\\
\bottomrule
\end{longtable}
Kolejność zakupów minimalizująca ryzyko: najpierw etap 0 (FP2+CDSR $=$ najdłuższe/najbardziej
niepewne pozycje), ogniwa dopiero po zaliczeniu CAN, UPS po potwierdzeniu modelu z szyną 48~V.

\section{Kosztorys zbiorczy (PL, lipiec 2026, $\pm$30\%)}
\begin{longtable}{@{}p{0.62\textwidth}r@{}}
\toprule
\textbf{MODUŁ A} & \textbf{zł}\\
\midrule
-X1 Type 2 z blokadą + odczep -X3 (7 żył) z wtykiem & 400--900\\
tor mocy: -Q0, -F1..3, szyny 6 mm$^2$, -K1, -X2 & 500--900\\
łańcuch bezpieczeństwa: E-STOP (2-torowy), -RB, -K10, -K4, -A2, krańcówka & 350--600\\
-B1 CDSR/RCM + generator -A3 (DAC, AMC1301, MOSFET) & 350--800\\
serwa DS3218 $\times$2 + rama + sprzęgła & 150--300\\
termowizja: MLX90640 + 2$\times$MLX90614 + DS18B20 $\times$8 & 400--600\\
metering: ATM90E36A + 3$\times$CT 40 A/1 V + dzielniki + -K30..33 & 200--400\\
logika: ESP32-S3, MCP23017, ULN2803, opto, RTC, SD, transceiver CAN, -G1, -F4/-F5 & 300--500\\
HMI DWIN 7$''$ + LED + przyciski & 300--500\\
obudowa IP54, przegrody, PCB, złącza, drobnica & 900--1600\\
\midrule
\textbf{Moduł A razem} & \textbf{3850--7100}\\
\midrule
\textbf{SZAFA B (wariant rekomendowany: B-b1 + W2)} & \\
\midrule
7$\times$ Flatpack2 48/2000 HE używane (6+1 zapas) + gniazda/półka & 1150--2400\\
-K11..-K13, rozdział AC, pole -X6, okablowanie AC & 350--600\\
szyna DC (płaskownik), -F13.1/2, -Q1, -KB, przewody 50 mm$^2$ & 600--1100\\
pakiet 16S1P: 16$\times$LFP 105 Ah kl. B + BMS JK + kompresja/busbary & 2740--4100\\
falownik W2: używany UPS on-line 5--6 kVA (48 V) + rozdzielniczka wyspy (B16+RCD+przełącznik) & 800--2000\\
szafa na kółkach, wentylatory (-M3/4, -K2/-K3), -X5, drobnica & 500--900\\
\midrule
\textbf{Szafa B razem} & \textbf{6140--11100}\\
\textbf{STACJA RAZEM} & \textbf{$\sim$9990--18200}\\
\quad rozbudowa 2P (+16 ogniw) & +2240--3200\\
\quad alternatywa W3 zamiast W2 (falownik własny LF) & $-$0--$+$700\\
\bottomrule
\end{longtable}
Plan pomostowy (niezależny): UT595 + UT210E + UT522 $\approx$ 2,8--3,4 tys. zł --- przyrządy te są
częścią systemu (pomiary protokołowe), zakup pozostaje zasadny od zaraz.

\section{Rejestr ryzyk (skonsolidowany R1--R24)}
\begin{longtable}{@{}p{0.05\textwidth}p{0.55\textwidth}p{0.34\textwidth}@{}}
\toprule
\textbf{Nr} & \textbf{Ryzyko} & \textbf{Mitygacja}\\
\midrule
R1 & PM701E na odczepie: poprawność pre-testów, prąd gniazd bananowych (brak instrukcji!) & test w etapie 0; pozyskać instrukcję\\
R2 & mechanika serw: swoboda pokręteł, trwałość & etap 0 na sprzęcie; fallback: własny front-end stanów CP\\
R3 & dostępność czujnika CDSR/RCM14 & zamówić najwcześniej; alternatywy WA RCM14\\
R4 & cieplne (teraz tylko straty 0,4--0,9 kW) & wentylatory + termiki w łańcuchu\\
R5 & generator upływu: potencjał na PE & ogranicznik sprzętowy, -K30, procedura\\
R6 & wzorcowanie bloków własnych & protokoły na UT595/UT522; wzorcowanie opcją\\
R7 & zgodność formalna (61010/61851/SEP) & przegląd przed etapem 2; uprawnienia E/D\\
R8 & spójność CP$\leftrightarrow$pobór & limity w firmware + pomiar -U2\\
R9--R10 & protokół CAN FP2 nieoficjalny; dynamika przy skokach & etap 0(c); rampy; fallback R4850G2\\
R11 & koordynacja SOC/BMS w soaku & bramka SOC; rozładowanie równoległe\\
R12 & formalności oddawania (dotyczy tylko upgrade W1) & start: wyspa bez zgłoszeń\\
R13 & duże prądy DC 48 V & przekroje, momenty, mega, przegląd\\
R14--R15 & jakość falownika DIY (W3); N--PE/RCD wyspy & procedury testowe; odbiór E/D (dot. też W2)\\
R20 & jakość ogniw używanych/kl. B & selekcja IR/pojemność, zakup z pomiarami, zapas\\
R21 & chemia NMC (gdyby moduły EV) & rekomendacja: LFP\\
R22 & egzemplarz UPS (stan, szyna 48 V, dokumentacja) & weryfikacja modelu przed zakupem, test wstępny\\
R23 & czas pracy własnej przekroczy plan & etapowanie, bufor czasowy, przepisy krok-po-kroku\\
R24 & 1P za mały na pełny soak 3F & bramka SOC $\le$25\%, rozładowanie równoległe, rozbudowa 2P\\
\bottomrule
\end{longtable}

\section{Decyzje projektowe --- status}
\begin{longtable}{@{}p{0.10\textwidth}p{0.64\textwidth}p{0.18\textwidth}@{}}
\toprule
\textbf{Nr} & \textbf{Decyzja} & \textbf{Status}\\
\midrule
Z1--Z6 & topologia odczepu PM701E; dwa moduły; mapa gniazd; pomiary protokołowe UT595/UT522;
etapowanie; budżet kierunkowy & zatwierdzone\\
Z7--Z10 & kierunek rekuperacyjny (PFC$\to$DC$\to$magazyn), parametry, tryb bez eksportu, budżet delta & zatwierdzone co do zasady\\
Z11--Z12 & falownik: warianty W1--W3; praca wyłącznie wyspowa do czasu W1 & W2 rekomendowany\\
Z13--Z15 & (wariant termiczny v4) & odrzucone przez zamawiającego\\
Z16 & kierunek v5/v6: budżetowy magazyn elektryczny (FP2 + pakiet DIY + W2) & \textbf{do zatwierdzenia}\\
Z17 & ogniwa: LFP kl. B 16S1P $\to$ 2P (alternatywy: moduły EV, rack) & \textbf{do zatwierdzenia}\\
Z18 & falownik: W2 (UPS) vs W3 (własny LF) & \textbf{do zatwierdzenia}\\
Z19 & budżet stacji $\sim$10--18,2 tys. zł & \textbf{do zatwierdzenia}\\
\bottomrule
\end{longtable}

\section{Testy odbiorcze i walidacja końcowa}
\begin{enumerate}
\item \textbf{Pakiet:} pojemność zmierzona $\ge$90\% deklarowanej; rozrzut ogniw pod obciążeniem
$\le$50~mV; zadziałanie progów BMS i -KB.
\item \textbf{Łańcuch E2:} kolejno wymuszone: E-STOP, termik, krańcówka, zanik wentylatorów,
zanik impulsów watchdoga, zanik 24~V --- każde musi otworzyć -K1; powrót tylko przez START;
E-STOP musi też gasić CP (-K4) --- DUT otwiera stycznik $\le$1~s.
\item \textbf{Wyspa:} RCD 30~mA (UT595: I$_{\Delta n}$, czas), praca 1--4~kW, udar silnikowy,
1~h soak termiczny.
\item \textbf{Tor pomiarowy:} zgodność -U2 z UT210E/UT890C ($\pm$2\%); CP duty vs deklaracje
pokręteł; RDC-DD na ładowarce referencyjnej (Q11): prąd $\le$6~mA, czas $<$ limitu.
\item \textbf{Test kompletny ładowarki referencyjnej} (P i Q11): pełna sekwencja 0--10 bez
interwencji poza wpisami UT595/UT522; raport wygenerowany; energia zmagazynowana zgodna z -U2
w granicach sprawności (85--92\%).
\end{enumerate}

\section{Dokumenty powiązane i źródła}
\textbf{Dokumenty projektu (folder \texttt{custom\_device}):} schemat elektryczny v5 (EPLAN,
E1--E3) + \emph{objaśnienie i przewodnik uruchomienia}; wizualizacja 3D v3 (HTML); koncepcje
v1--v5 (historia decyzji); \texttt{custom\_device\_inf.md} (kontekst).
\textbf{Źródła:} IEC 61851-1 (interfejs CP/PP); IEC 62955 (RDC-DD); IEC 61010 (bezpieczeństwo
przyrządów); PN-HD 60364 (instalacje, wyspa, N--PE); PN-EN 50549-1/NC RfG (dot. tylko przyszłego
W1); noty: EVE LF105, smart BMS 16S (JK), Eltek Flatpack2 48/2000 HE + protokół CAN (dokumentacja
społecznościowa openinverter/endless-sphere), LEM CDSR, ATM90E36A, MLX90640/14, EG8010 (wariant W3);
,,Przewodnik pomiarów elektrycznych stacji ładowania'' (Warszawa 2024).

\end{document}
